%==============================================================================
% jarvis-listings.tex
% Code-listing setup for the Jarvis-HEP manual.
%
% Provides:
%   - a Jarvis color palette
%   - a custom YAML language definition (listings has none built in)
%   - styles:  yaml   (task-card YAML)
%              bash   (shell / CLI sessions)
%              plain  (generic / SLHA / output dumps)
%   - convenience inline/float macros
%
% Long real-world cards should be \lstinputlisting'd from ../listings/ so they
% stay byte-true with jarvishep/project_template.  Inside Tufte's narrow text
% column, wrap wide listings in a {fullwidth} environment.
%==============================================================================

\usepackage{listings}

% --- Palette ------------------------------------------------------------------
\definecolor{jBg}{HTML}{F7F7F4}      % listing background
\definecolor{jRule}{HTML}{D8D8D0}    % frame rule
\definecolor{jKey}{HTML}{1F4E79}     % YAML keys / keywords (deep blue)
\definecolor{jStr}{HTML}{8A5A00}     % strings (amber)
\definecolor{jComment}{HTML}{6E7B6E} % comments (muted green-grey)
\definecolor{jNum}{HTML}{7A1F8B}     % numbers (purple)
\definecolor{jMarker}{HTML}{B3261E}  % &J/ @tokens / ${...} (Jarvis red)
\definecolor{jLineNo}{HTML}{A0A0A0}
\definecolor{jLineNum}{HTML}{C7CCD2}  % very light grey for YAML line numbers
\definecolor{jGuide}{HTML}{E3E6EA}    % very light grey for YAML indentation guides

% --- Semantic colours (used inline and in listings) ---------------------------
\definecolor{jExpr}{HTML}{C2185B}     % expressions  -> pink
\definecolor{jPath}{HTML}{1B5E20}     % file/dir paths -> dark green
\definecolor{jCmd}{HTML}{E8730C}      % Jarvis CLI commands -> Jarvis orange

% --- Inline semantic markup ---------------------------------------------------
\newcommand{\expr}[1]{\textcolor{jExpr}{\jtt{#1}}}   % an expression, e.g. \expr{x * Pi}
\newcommand{\cmd}[1]{\textcolor{jCmd}{\jtt{#1}}}     % a Jarvis CLI command, e.g. \cmd{Jarvis}

% --- A minimal YAML dialect ---------------------------------------------------
% Good enough to colour keys, strings, comments, booleans, and Jarvis markers
% without a full parser.
\lstdefinelanguage{jyaml}{
  morekeywords={true,false,null,True,False,None,yes,no},
  sensitive=true,
  morecomment=[l]{\#},
  morestring=[b]",
  morestring=[b]',
  alsoletter={-,_,.},
}

% --- Shared base style --------------------------------------------------------
\lstdefinestyle{jbase}{
  backgroundcolor=\color{jBg},
  basicstyle=\ttfamily\footnotesize,
  breaklines=true,
  breakatwhitespace=false,
  postbreak=\mbox{\textcolor{jLineNo}{$\hookrightarrow$}\space},
  columns=fullflexible,
  keepspaces=true,
  showstringspaces=false,
  tabsize=2,
  frame=leftline,
  framerule=1.2pt,
  rulecolor=\color{jRule},
  xleftmargin=0.8em,
  framexleftmargin=0.8em,
  aboveskip=0.9\baselineskip,
  belowskip=0.6\baselineskip,
}

% --- YAML style ---------------------------------------------------------------
\lstdefinestyle{yaml}{
  style=jbase,
  language=jyaml,
  keywordstyle=\color{jNum}\bfseries,   % booleans/null
  stringstyle=\color{jStr},
  commentstyle=\color{jComment}\itshape,
  % Highlight Jarvis markers and tokens.
  moredelim=[is][\color{jMarker}\bfseries]{@@}{@@},
  % Very light line numbers in the left margin (outside the frame rule).
  numbers=left,
  numberstyle=\tiny\color{jLineNum},
  numbersep=9pt,
  xleftmargin=2.6em,
  framexleftmargin=2.6em,
}

% --- Bash / CLI style ---------------------------------------------------------
\lstdefinestyle{bash}{
  style=jbase,
  language=bash,
  keywordstyle=\color{jKey}\bfseries,
  commentstyle=\color{jComment}\itshape,
  stringstyle=\color{jStr},
  morekeywords={Jarvis,pip,python3,cd,latexmk},
}

% --- Plain dump style (SLHA blocks, output text) ------------------------------
\lstdefinestyle{plain}{
  style=jbase,
  language={},
  keywordstyle=\color{jKey},
  commentstyle=\color{jComment}\itshape,
}

% --- Inner style for the boxed shell environment (frame/bg come from the box) --
\lstdefinestyle{shellinner}{
  language=bash,
  basicstyle=\ttfamily\footnotesize,
  keywordstyle=\color{jKey}\bfseries,
  commentstyle=\color{jComment}\itshape,
  stringstyle=\color{jStr},
  morekeywords={pip,python3,cd,source,pyenv,latexmk,make},
  keywords=[2]{Jarvis,jplot,jopera},        % Jarvis-family executables
  keywordstyle=[2]\color{jCmd}\bfseries,     % ... in Jarvis orange
  breaklines=true,
  postbreak=\mbox{\textcolor{jLineNo}{$\hookrightarrow$}\space},
  columns=fullflexible, keepspaces=true, showstringspaces=false, tabsize=2,
  aboveskip=0pt, belowskip=0pt,   % no extra listing space inside the box
}

% --- YAML style for the "window" box (frame/background come from the box) ------
% Line numbers (very light), expression values in pink, &J paths in dark green.
\lstdefinestyle{yamlwindow}{
  language=jyaml,
  basicstyle=\ttfamily\footnotesize,
  keywordstyle=\color{jNum}\bfseries,
  commentstyle=\color{jComment}\itshape,
  morestring=[b]", stringstyle=\color{jStr},
  moredelim=[s][\color{jExpr}]{expression: "}{"},
  moredelim=[l][\color{jPath}]{"\&J},
  numbers=left, numberstyle=\tiny\color{jLineNum}, numbersep=12pt,
  xleftmargin=2.4em,
  breaklines=true, columns=fullflexible, keepspaces=true, showstringspaces=false, tabsize=2,
  aboveskip=0pt, belowskip=0pt,
}

\lstset{style=jbase}

% --- Convenience environments / macros ----------------------------------------
% \begin{yamlcode} ... \end{yamlcode}
\lstnewenvironment{yamlcode}{\lstset{style=yaml}}{}
% \begin{bashcode} ... \end{bashcode}
\lstnewenvironment{bashcode}{\lstset{style=bash}}{}
% \begin{plaincode} ... \end{plaincode}
\lstnewenvironment{plaincode}{\lstset{style=plain}}{}

% File include shortcuts (keep wide ones inside {fullwidth}).
\newcommand{\yamlfile}[1]{\lstinputlisting[style=yaml]{#1}}
\newcommand{\bashfile}[1]{\lstinputlisting[style=bash]{#1}}

% Inline monospace shortcut already provided as \code in the preamble.
